\section{Limitations}
\label{sec:limitations}
% =====================================================================

A limitation intrinsic to the setting is that a single stationary snapshot cannot
identify a dynamical system. Theorem~\ref{thm:identifiability} makes this precise:
for $C\ge3$ the stationarity constraint determines $f_\theta$ only on the range
$\mathcal{R}=\vz^\star(\Omega)$, a Lebesgue-null set, so counterfactuals that
displace $\vz$ off $\mathcal{R}$ are governed by the architecture rather than the
data. Reported diffusion lengths accordingly characterize the fitted operator
rather than physical morphogen transport; Section~\ref{sec:results} shows
they are moreover close to their initialization, which is why we make no
claim about signaling and why the operator is named for its mathematical
class rather than for a biological one. Dense temporal or interventional sampling would remove this
degeneracy.

We note that the predicted fields are smoother than the measurements, and that on
Moran's $I$ error NRDO trails the baselines it leads on the other metrics. This is
a known consequence of diffusion-based architectures
\citep{grand2021,gread2022} and bounds the applications we would recommend:
prediction at unmeasured coordinates is well supported, delineation of sharp
boundaries less so. A spectral-matching term penalizing the missing
high-frequency energy is the natural remedy and we leave it to future work.

The benchmark's statistical resolution is limited by the number of independent
specimens rather than by the number of sections. Twelve of \MSSections{} sections
come from one brain, leaving \MSSpecPairs{} independent specimens, at which the
smallest attainable Wilcoxon $p$ is $\MSSpecMinP$. Against the continuous
baselines this is enough: NRDO wins on every specimen, which is the strongest
result the design admits. Against the graph baselines it is not
($d_z = \MSSpecGraphDz$, NRDO ahead on \MSSpecWins{} of \MSSpecPairs{}); the
section-level separation there rests substantially on correlated sections of one
brain, and additional independent tissues --- not additional sections --- are
what would settle it. We regard this as the single most informative experiment
left undone.

Transfer remains the weakest setting. NRDO retains a substantial fraction of
in-domain performance zero-shot, but the STAGATE-style baseline transfers better,
and until $\mD_\theta$ is conditioned on tissue-level covariates the operator is
best regarded as a within-section model. Finally, the \emph{-style} baselines
implement each method's architectural core with a common read-out head and are
not reproductions of published pipelines; the benchmark uses a reduced epoch
budget and subsampled sections so that \MSRuns{} runs were feasible on CPU, which
lowers all models' absolute scores; and scope is limited to single 2-D sections.

\iffullpaper\else\label{endofmain}\fi
